\documentclass[UTF8,fontset=none,linespread=1]{ctexart}
% Explicit font families avoid mixed CJK fallbacks in PDF viewers.
\setmainfont{Helvetica Neue}[BoldFont={Helvetica Neue Bold},ItalicFont={Helvetica Neue Italic}]
\setsansfont{Helvetica Neue}
\setCJKmainfont{PingFang SC}[BoldFont={PingFang SC Semibold}]
\setCJKsansfont{PingFang SC}[BoldFont={PingFang SC Semibold}]
\setCJKmonofont{PingFang SC}
\xeCJKsetup{PunctStyle=kaiming,CJKecglue={\hskip .15em plus .05em minus .05em}}
\linespread{1}\selectfont
\hyphenpenalty=10000
\exhyphenpenalty=10000
\usepackage[paperwidth=280mm,paperheight=157.5mm,margin=0mm]{geometry}
\usepackage{amsmath,amssymb,bm,mathtools,tikz,hyperref}
\usetikzlibrary{arrows.meta,calc,decorations.pathmorphing,positioning}
\hypersetup{colorlinks=true,urlcolor=teal,pdftitle={末态 DiFF 探测四费米子相互作用：研究 proposal},pdfauthor={Research proposal},pdfsubject={EIC; helicity interference; collinear factorization}}
\pagestyle{empty}
\setlength{\parindent}{0pt}
\definecolor{ink}{HTML}{182C40}
\definecolor{blue}{HTML}{285D8F}
\definecolor{teal}{HTML}{167D8D}
\definecolor{orange}{HTML}{D26A39}
\definecolor{gray}{HTML}{657585}
\definecolor{pale}{HTML}{EEF4F7}
\definecolor{line}{HTML}{D8E1E7}
\newcommand{\bodyfont}{\fontsize{15}{21}\selectfont}
\newcommand{\smallfont}{\fontsize{11.5}{16.5}\selectfont}
\newcommand{\tinyfont}{\fontsize{9}{12}\selectfont}
\newcommand{\HT}{H_1^{\sphericalangle}}
\newcommand{\dd}{\mathrm d}
\newcommand{\src}[2]{\href{https://arxiv.org/abs/#1}{#2}}
\newcommand{\slide}[3]{\clearpage\null\begin{tikzpicture}[remember picture,overlay,x=1cm,y=1cm,shift={(current page.south west)},every node/.style={inner sep=0pt,text=ink},>=Latex]
\fill[white] (0,0) rectangle (28,15.75);
\draw[teal,line width=1.5pt] (1,14.72)--(2,14.72);
\node[anchor=west,font=\smallfont,text=gray] at (2.22,14.72) {RESEARCH PROPOSAL\quad /\quad EIC FOUR-FERMION};
\node[anchor=north west,font=\fontsize{25}{31}\selectfont\bfseries] at (1,14.03) {#2};
\draw[line] (1,.82)--(27,.82);
\node[anchor=west,font=\tinyfont,text=gray] at (1,.44) {#3};
\node[anchor=east,font=\tinyfont,text=gray] at (27,.44) {#1\,/\,10};}
\newcommand{\finishslide}{\end{tikzpicture}}
\newcommand{\txt}[4]{\node[anchor=north west,text width=#3cm,font=\bodyfont,align=left] at (#1,#2) {#4};}
\newcommand{\sm}[4]{\node[anchor=north west,text width=#3cm,font=\smallfont,align=left] at (#1,#2) {#4};}
\tikzset{ferm/.style={-{Latex[length=2.5mm]},line width=1.3pt,draw=blue},had/.style={-{Latex[length=2.5mm]},line width=1.3pt,draw=orange},spin/.style={-{Latex[length=2.7mm]},line width=1.7pt,draw=teal},flow/.style={-{Latex[length=2.5mm]},line width=1.1pt,draw=gray}}
\begin{document}

\slide{1}{末态 DiFF 探测四费米子相互作用}{核心问题与研究假设}
\txt{1}{12.65}{25.7}{\textbf{目标：}用末态双强子面读取 quark transverse spin，研究 scalar/tensor 四费米算符与 SM 的线性干涉。电子横极化保留，proton 可不极化。}
\node[anchor=west,font=\fontsize{24}{30}\selectfont,text=blue] at (1.05,10.38) {$e_T^-p\longrightarrow e^-\,[\pi^+\pi^-]_{\rm current}+X$};
% Cover schematic: the same hard/fragmentation topology as slide 4.
\begin{scope}
\coordinate (C1H) at (7.3,5.9);
\coordinate (C1F) at (14.4,5.9);
\coordinate (C1Pone) at (19.4,8.0);
\coordinate (C1Ptwo) at (19.4,3.8);
% Reusing these endpoints keeps the pion arrows on the filled plane edges.
\fill[orange!10] (C1F)--(C1Pone)--(C1Ptwo)--cycle;
\draw[orange!30,line width=.7pt] (C1Pone)--(C1Ptwo);

\draw[ferm,line width=1.65pt,shorten >=3.7mm] (1.65,8.3)--(C1H);
\draw[ferm,line width=1.65pt] (C1H)--(10.5,8.3);
\node[anchor=south west,font=\fontsize{20}{26}\selectfont,text=blue] at (1.65,8.62) {$e_T^-$};
\node[anchor=south,font=\fontsize{20}{26}\selectfont,text=blue] at (10.50,8.64) {$e^-$};
\draw[spin,line width=1.75pt] (3.51,7.51)--(4.01,8.69);
\node[font=\bodyfont,text=teal,anchor=south] at (4.10,8.98) {$S_{eT}$};

\draw[-{Latex[length=3mm]},draw=ink,line width=1.65pt,shorten >=3.7mm] (1.65,3.4)--(C1H);
\node[anchor=north west,font=\bodyfont,text=ink] at (1.65,2.99) {入射 quark};
\node[anchor=north west,font=\smallfont,text=gray] at (1.65,2.28) {来自非极化 proton};
\draw[-{Latex[length=3mm]},draw=ink,line width=1.65pt,shorten >=2.5mm] (C1H)--(C1F);
\node[font=\bodyfont,text=ink] at (10.90,5.25) {出射 quark $\ell$};
\draw[spin,line width=1.75pt] (11.70,5.9)--(11.70,7.28);
\node[font=\bodyfont,text=teal,anchor=south] at (11.70,7.56) {$s_{qT}$};
\filldraw[fill=pale,draw=ink!60,line width=1pt] (C1H) circle (.36);
\node[font=\fontsize{20}{25}\selectfont,text=ink] at (C1H) {$\mathcal H$};
\node[font=\smallfont,text=gray] at (8.70,4.25) {SM--EFT interference};

\draw[gray!80,densely dashed,-{Latex[length=2.3mm]},line width=1pt] (C1F)--(20.05,5.9);
\node[font=\bodyfont,text=gray,anchor=south] at (18.45,6.16) {$P_h$};
\draw[had,line width=1.8pt] (C1F)--(C1Pone);
\draw[had,line width=1.8pt] (C1F)--(C1Ptwo);
\node[font=\fontsize{20}{26}\selectfont,text=orange,anchor=west] at (19.67,8.02) {$\pi^+$};
\node[font=\fontsize{20}{26}\selectfont,text=orange,anchor=west] at (19.67,3.78) {$\pi^-$};
\node[font=\bodyfont,text=orange] at (17.10,8.96) {dihadron plane};
\draw[gray!65,densely dashed,-{Latex[length=2.2mm]},line width=1pt] (C1F)--(15.27,3.14);
\node[font=\smallfont,text=gray,anchor=north] at (15.28,2.81) {$X_{\rm frag}$};
\filldraw[fill=orange!15,draw=orange,line width=1pt] (C1F) circle (.24);

% The measured angle is defined in an end-on view, not by perspective.
\node[font=\smallfont,text=gray] at (24.08,8.98) {末态角度读出};
\begin{scope}[shift={(24.05,5.9)}]
\fill[orange!7] (0,0)--(.88,0) arc[start angle=0,end angle=80,radius=.88]--cycle;
\draw[line,line width=.7pt] (0,0) circle (1.65);
\draw[gray!65,-{Latex[length=2.1mm]},line width=.8pt] (-1.86,0)--(1.95,0);
\draw[gray!65,-{Latex[length=2.1mm]},line width=.8pt] (0,-1.80)--(0,1.97);
\node[font=\smallfont,text=gray,anchor=west] at (2.06,0) {$\hat x$};
\node[font=\smallfont,text=gray,anchor=south] at (0,2.07) {$\hat y$};
\draw[had,line width=1.75pt] (0,0)--({1.53*cos(80)},{1.53*sin(80)});
\node[font=\bodyfont,text=orange,anchor=west] at (.51,1.52) {$R_T$};
\draw[spin,line width=1.75pt] (0,0)--({1.5*cos(145)},{1.5*sin(145)});
\node[font=\smallfont,text=teal,anchor=east] at (-1.20,1.17) {$s_{qT}$};
\draw[orange,-{Latex[length=2mm]},line width=1pt] (.88,0) arc[start angle=0,end angle=80,radius=.88];
\node[font=\fontsize{19}{24}\selectfont,text=orange] at (1.06,.78) {$\phi_R$};
\fill[white] (0,0) circle (.11);
\draw[ink,line width=.75pt] (0,0) circle (.075);
\fill[ink] (0,0) circle (.023);
\node[font=\smallfont,anchor=north west] at (.16,-.15) {$\hat z$};
\end{scope}
\node[font=\smallfont,text=gray] at (24.05,3.54) {从 $+\hat z$ 侧看向原点};
\end{scope}
\finishslide

\slide{2}{已有可参考的文章：四费米算符}{[1] Wang, Wen, Xing, Yan, PRD 109, 095025 (2024), Eq. (11)；示意图依据算符结构绘制}
\node[anchor=west,font=\bodyfont,text=gray] at (1,12.42) {$e_T^-p_T\to e^-X$：电子横极化 $+$ proton transversity};
\draw[line] (13.65,11.60)--(13.65,2.35);
\node[anchor=west,font=\fontsize{17}{23}\selectfont] at (1,11.26) {$\mathcal L=\mathcal L_{\rm SM}+\sum_i C_i\mathcal O_i/\Lambda^2+{\rm h.c.}$};
\node[anchor=west,font=\fontsize{18}{25}\selectfont,text=blue] at (1,9.65) {$\mathcal O_{ledq}=(\bar L^j e)(\bar d Q^j)$};
\node[anchor=west,font=\fontsize{14}{20}\selectfont] at (1.15,8.80) {$e$--$d$ scalar：\ $(\bar e_L e_R)(\bar d_R d_L)$};
\node[anchor=west,font=\fontsize{18}{25}\selectfont,text=blue] at (1,7.48) {$\mathcal O_{lequ}^{(1)}=(\bar L^j e)\epsilon_{jk}(\bar Q^k u)$};
\node[anchor=west,font=\fontsize{14}{20}\selectfont] at (1.15,6.60) {$e$--$u$ scalar：\ $-(\bar e_L e_R)(\bar u_L u_R)$};
\node[anchor=west,font=\fontsize{17.5}{24}\selectfont,text=blue] at (1,5.23) {$\mathcal O_{lequ}^{(3)}=(\bar L^j\sigma_{\mu\nu}e)\epsilon_{jk}(\bar Q^k\sigma^{\mu\nu}u)$};
\node[anchor=west,font=\fontsize{14}{20}\selectfont] at (1.15,4.34) {$e$--$u$ tensor：\ $-(\bar e_L\sigma_{\mu\nu}e_R)(\bar u_L\sigma^{\mu\nu}u_R)$};
\node[anchor=west,font=\smallfont,text=gray] at (1.1,2.83) {各算符下方列 neutral-current 部分，另加 h.c.};

\node[anchor=west,font=\bodyfont\bfseries,text=blue] at (14.6,11.65) {SM：两条线保持 chirality};
\draw[ferm,line width=1.4pt] (16.2,10.70)--(24.3,10.70);
\draw[-{Latex[length=2.5mm]},draw=ink,line width=1.4pt] (16.2,9.12)--(24.3,9.12);
\draw[decorate,decoration={snake,amplitude=.6mm,segment length=2.3mm},gray,line width=1pt] (20.25,10.70)--(20.25,9.12);
\fill[blue] (20.25,10.70) circle (.06);
\fill[ink] (20.25,9.12) circle (.06);
\node[font=\bodyfont,text=blue,anchor=east] at (15.97,10.70) {$e_L$};
\node[font=\bodyfont,text=blue,anchor=west] at (24.55,10.70) {$e_L$};
\node[font=\bodyfont,anchor=east] at (15.97,9.12) {$q_{R/L}$};
\node[font=\bodyfont,anchor=west] at (24.55,9.12) {$q_{R/L}$};
\node[font=\smallfont,text=gray,anchor=west] at (20.68,9.91) {$\gamma/Z$};

\node[anchor=west,font=\bodyfont\bfseries,text=orange] at (14.6,7.79) {scalar / tensor：两条线翻转 chirality};
\coordinate (O2V) at (20.25,5.68);
\draw[ferm,line width=1.4pt,shorten >=1.4mm] (16.2,6.82)--(O2V);
\draw[ferm,line width=1.4pt] (O2V)--(24.3,6.82);
\draw[-{Latex[length=2.5mm]},draw=ink,line width=1.4pt,shorten >=1.4mm] (16.2,4.54)--(O2V);
\draw[-{Latex[length=2.5mm]},draw=ink,line width=1.4pt] (O2V)--(24.3,4.54);
\fill[orange] (O2V) circle (.13);
\node[font=\bodyfont,text=orange,anchor=south] at (20.25,6.13) {$\mathcal O_i$};
\node[font=\bodyfont,text=blue,anchor=east] at (15.97,6.82) {$e_R$};
\node[font=\bodyfont,text=blue,anchor=west] at (24.55,6.82) {$e_L$};
\node[font=\bodyfont,anchor=east] at (15.97,4.54) {$q_{L/R}$};
\node[font=\bodyfont,anchor=west] at (24.55,4.54) {$q_{R/L}$};
\node[font=\fontsize{14}{20}\selectfont] at (20.5,3.38) {$ledq:\ d_L\to d_R\qquad lequ^{(1,3)}:\ u_R\to u_L$};
\finishslide

\slide{3}{原文的角投影：为什么有 $2\phi$，也有常数项？}{[1] Eq. (1), Eqs. (12)--(16)；共同横平面中的角度示意，取两束横自旋平行}
\node[anchor=west,font=\fontsize{14}{20}\selectfont\bfseries,text=teal] at (1.2,12.35) {inclusive DIS：$e_T^-(k)+p_T(P)\to e^-(k')+X$};
\begin{scope}[shift={(5.30,8.22)}]
\draw[line,line width=.7pt] (0,0) circle (2.45);
\fill[orange!8] (0,0)--(1.30,0) arc[start angle=0,end angle=65,radius=1.30]--cycle;
\draw[gray!50,dashed,line width=.7pt] (-2.7,0)--(3.05,0);
\draw[ferm,line width=1.65pt] (0,0)--(3.05,0);
\node[font=\bodyfont,text=blue,anchor=west] at (3.22,0) {$k'_T$};
\draw[spin,line width=1.8pt] (0,0)--({2.8*cos(65)},{2.8*sin(65)});
\node[font=\fontsize{16}{22}\selectfont,text=teal,anchor=south] at (.85,2.96) {$S_{eT}\parallel S_{pT}$};
\draw[orange,-{Latex[length=2.2mm]},line width=1.15pt] (1.30,0) arc[start angle=0,end angle=65,radius=1.30];
\node[font=\fontsize{25}{31}\selectfont,text=orange] at (1.63,1.05) {$\phi$};
\fill[ink] (0,0) circle (.055);
\end{scope}
\node[font=\bodyfont,text=gray] at (5.35,5.12) {lab frame：沿共同束流轴看};
\node[font=\fontsize{22}{29}\selectfont,text=teal] at (5.7,3.96) {$\phi_1=\phi_2=\phi$};
\node[font=\smallfont,text=gray] at (5.7,3.00) {$\phi_1=\angle(k'_T,S_{eT}),\quad\phi_2=\angle(k'_T,S_{pT})$};

\fill[pale,rounded corners=3pt] (12.2,7.39) rectangle (27,12.12);
\node[anchor=west,font=\bodyfont\bfseries,text=blue] at (12.7,11.50) {$ledq$：相反 chiral projectors $\Rightarrow$ 和角};
\node[anchor=west,font=\fontsize{21}{28}\selectfont] at (12.7,10.32) {$\Phi_+=\phi_1+\phi_2=2\phi$};
\node[anchor=west,font=\fontsize{18}{25}\selectfont] at (12.7,9.20) {$\operatorname{Re}[Ce^{-2i\phi}]=\operatorname{Re}C\cos2\phi+\operatorname{Im}C\sin2\phi$};
\node[anchor=west,font=\bodyfont,text=orange] at (12.7,8.10) {$w=\cos2\phi,\ \sin2\phi$};

\fill[orange!5,rounded corners=3pt] (12.2,2.37) rectangle (27,6.82);
\node[anchor=west,font=\bodyfont\bfseries,text=blue] at (12.7,6.15) {$lequ^{(1,3)}$：相同 chiral projectors $\Rightarrow$ 差角};
\node[anchor=west,font=\fontsize{21}{28}\selectfont] at (12.7,5.01) {$\Phi_-=\phi_1-\phi_2=0$};
\node[anchor=west,font=\fontsize{20}{27}\selectfont] at (12.7,3.96) {$\operatorname{Re}[Ce^{-i\Phi_-}]=\operatorname{Re}C\qquad w=1$};
\node[anchor=west,font=\smallfont,text=gray] at (12.7,2.94) {scalar / tensor 改变 $y$ 依赖；这里的差角都不随 $\phi$ 变化。};
\finishslide

\slide{4}{把 quark 自旋分析从初态移到末态}{电子横极化保留；quark 的 chiral-odd 自旋矩阵元由初态 transversity 改为末态 DiFF}
\node[font=\fontsize{20}{27}\selectfont\bfseries,text=blue] at (7.0,12.18) {原方案};
\node[font=\fontsize{20}{27}\selectfont\bfseries,text=orange] at (21.0,12.18) {本方案};
\draw[line] (14,11.70)--(14,4.3);
\foreach \offset in {0,14}{
\begin{scope}[shift={(\offset,0)}]
\draw[ferm,line width=1.45pt] (2.0,9.97)--(11.7,9.97);
\draw[-{Latex[length=2.5mm]},draw=ink,line width=1.45pt] (2.0,7.30)--(11.7,7.30);
\filldraw[fill=pale,draw=ink!30,line width=.85pt,rounded corners=4pt] (6.0,6.67) rectangle (8.3,10.60);
\node[font=\fontsize{23}{30}\selectfont] at (7.15,8.64) {$\mathcal H$};
\node[font=\smallfont,text=blue,anchor=east] at (1.80,9.97) {$e$};
\node[font=\smallfont,anchor=east] at (1.80,7.30) {$q$};
\draw[teal,line width=3pt] (4.0,9.64)--(4.0,10.30);
\node[font=\bodyfont,text=teal,anchor=south] at (4.0,10.73) {$S_{eT}$};
\node[font=\smallfont,text=teal] at (4.0,9.18) {chiral-odd};
\end{scope}}
\draw[orange,line width=3pt] (4.0,6.97)--(4.0,7.63);
\node[font=\fontsize{22}{28}\selectfont,text=orange] at (4.0,5.84) {$h_1^q$};
\node[font=\smallfont,text=gray] at (4.0,4.99) {初态 transversity};
\draw[orange,line width=3pt] (24.0,6.97)--(24.0,7.63);
\node[font=\fontsize{22}{28}\selectfont,text=orange] at (24.0,5.84) {$H_{1,q}^{\sphericalangle}$};
\node[font=\smallfont,text=gray] at (24.0,4.99) {末态 DiFF};
\node[font=\smallfont,text=gray] at (17.95,6.24) {$f_1^q$};
\node[font=\fontsize{22}{30}\selectfont] at (14,2.68) {$e\text{ line}:\ S_{eT}\qquad\qquad q\text{ line}:\ h_1^q\ \longrightarrow\ H_{1,q}^{\sphericalangle}$};
\finishslide


\slide{5}{双强子面怎样读取 quark transverse spin}{[2] Eqs. (2)--(4)；采用显式保留 $|R_T|/M_h$ 的 interference DiFF 记号}
% The hard-scattering part is a topology schematic, not momentum geometry.
% Pair-plane fill and pion arrows share exactly the same three coordinates.
\begin{scope}
\node[anchor=west,font=\smallfont,text=blue] at (1.05,12.1) {硬散射（示意）};
\draw[blue!45,line width=.65pt] (1.05,11.72)--(4.55,11.72);
\node[anchor=west,font=\smallfont,text=orange] at (5.9,12.1) {quark fragmentation};
\draw[orange!50,line width=.65pt] (5.9,11.72)--(9.1,11.72);
\node[font=\smallfont,text=gray] at (11.65,12.1) {横向截面};

\coordinate (D4H) at (3.1,9.3);
\coordinate (D4F) at (6.5,9.3);
\coordinate (D4Pone) at (8.95,10.62);
\coordinate (D4Ptwo) at (8.95,7.98);
\coordinate (D4PairAxis) at (9.25,9.3);
\fill[orange!10] (D4F)--(D4Pone)--(D4Ptwo)--cycle;
\draw[orange!30,line width=.55pt] (D4Pone)--(D4Ptwo);

% Both incoming and outgoing electron/quark legs meet the hard blob.
\draw[ferm,line width=1.3pt,shorten >=3mm] (1.2,11.0)--(D4H);
\draw[ferm,line width=1.3pt] (D4H)--(4.9,11.0);
\node[anchor=south,font=\smallfont,text=blue] at (1.55,11.14) {$e^-(k)$};
\node[anchor=south,font=\smallfont,text=blue] at (4.63,11.14) {$e^-(k')$};
\draw[-{Latex[length=2.5mm]},draw=ink,line width=1.35pt,shorten >=3mm] (1.2,7.55)--(D4H);
\node[anchor=north west,font=\smallfont,text=ink] at (1.04,7.29) {入射 quark $p$};
\draw[-{Latex[length=2.5mm]},draw=ink,line width=1.35pt,shorten >=2mm] (D4H)--(D4F);
\node[font=\smallfont,text=ink] at (4.82,8.85) {出射 quark $\ell$};
\filldraw[fill=pale,draw=ink!60,line width=.85pt] (D4H) circle (.29);
\node[font=\fontsize{13}{17}\selectfont,text=ink] at (D4H) {$\mathcal H$};
\node[anchor=west,font=\fontsize{9.3}{13}\selectfont,text=gray] at (3.42,8.07) {SM--EFT interference};

% The two pion arrows lie exactly on the sides of the filled pair plane.
\draw[gray!80,densely dashed,-{Latex[length=1.8mm]},line width=.8pt] (D4F)--(D4PairAxis);
\node[anchor=south,font=\smallfont,text=gray] at (8.25,9.48) {$P_h$};
\draw[had,line width=1.4pt] (D4F)--(D4Pone);
\draw[had,line width=1.4pt] (D4F)--(D4Ptwo);
\node[anchor=south,font=\smallfont,text=orange] at (8.63,10.83) {$p_1\;(\pi^+)$};
\node[anchor=north,font=\smallfont,text=orange] at (8.63,7.74) {$p_2\;(\pi^-)$};
\draw[gray!65,densely dashed,-{Latex[length=1.8mm]},line width=.8pt] (D4F)--(7.03,7.55);
\node[anchor=north,font=\smallfont,text=gray] at (7.03,7.27) {$X_{\rm frag}$};
\filldraw[fill=orange!15,draw=orange,line width=.9pt] (D4F) circle (.19);
\node[font=\fontsize{10}{14}\selectfont,text=gray] at (2.72,6.52) {初态 quark 来自 proton};
\node[font=\fontsize{10}{14}\selectfont,text=gray] at (7.75,6.52) {$P_h\parallel\ell$，两强子来自同一 quark};

% End-on view from +z: x right, y up, z out of the page.
\begin{scope}[shift={(11.65,9.3)}]
\fill[orange!7] (0,0)--(.82,0) arc[start angle=0,end angle=80,radius=.82]--cycle;
\draw[line,line width=.6pt] (0,0) circle (1.5);
\draw[gray!65,-{Latex[length=1.8mm]},line width=.7pt] (-1.68,0)--(1.75,0);
\draw[gray!65,-{Latex[length=1.8mm]},line width=.7pt] (0,-1.63)--(0,1.76);
\node[font=\smallfont,text=gray,anchor=west] at (1.82,0) {$\hat x$};
\node[font=\smallfont,text=gray,anchor=south] at (0,1.84) {$\hat y$};
\draw[had,line width=1.6pt] (0,0)--({1.42*cos(80)},{1.42*sin(80)});
\node[font=\bodyfont,text=orange,anchor=west] at (.44,1.42) {$R_T$};
\draw[spin,line width=1.6pt] (0,0)--({1.38*cos(145)},{1.38*sin(145)});
\node[font=\bodyfont,text=teal,anchor=east] at (-1.14,1.09) {$s_{qT}$};
\draw[orange,-{Latex[length=1.7mm]},line width=.9pt] (.82,0) arc[start angle=0,end angle=80,radius=.82];
\node[font=\bodyfont,text=orange] at (1.0,.72) {$\phi_R$};
\fill[white] (0,0) circle (.095);
\draw[ink,line width=.65pt] (0,0) circle (.065);
\fill[ink] (0,0) circle (.018);
\node[font=\smallfont,anchor=north west] at (.13,-.12) {$\hat z$};
\end{scope}
\node[font=\fontsize{10}{14}\selectfont,text=gray] at (11.65,7.11) {从 $+\hat z$ 侧看向原点};
\end{scope}
\node[anchor=west,font=\fontsize{23}{31}\selectfont] at (15.2,11.93) {$P_h=p_1+p_2,\quad M_h^2=P_h^2$};
\node[anchor=west,font=\fontsize{23}{31}\selectfont] at (15.2,10.32) {$R=\frac12(p_1-p_2)$};
\node[anchor=west,font=\fontsize{22}{30}\selectfont] at (15.2,8.73) {$\bm R_T=|R_T|(\cos\phi_R\,\hat x+\sin\phi_R\,\hat y)$};
\node[anchor=west,font=\fontsize{16}{22}\selectfont,text=gray] at (15.2,7.17) {$\hat z\parallel\bm P_h,\quad\hat y\parallel\bm k\times\bm k',\quad\hat x=\hat y\times\hat z$};
\node[anchor=west,font=\fontsize{21}{29}\selectfont] at (1.05,4.55) {$\mathcal F_q=D_1^q+\dfrac{|R_T|}{M_h}H_{1,q}^{\sphericalangle}\bigl(s_{qy}\cos\phi_R-s_{qx}\sin\phi_R\bigr)$};
\node[anchor=west,font=\bodyfont,text=gray] at (1.1,2.73) {$D_1$：双强子产额\qquad $H_1^{\sphericalangle}$：quark transverse spin 的末态分析};
\finishslide

\slide{6}{从这个课题学习 QCD 因子化}{[2] Eqs. (2)--(4)；PDF、DiFF 与 hard response 的因子化结构}
\node[font=\fontsize{24}{32}\selectfont] at (14,10.75) {$\dd\sigma\ \propto\ \sum_q\underbrace{f_1^q}_{\text{PDF}}\ \otimes\ \underbrace{W^q(\rho^e,C_i)}_{\text{hard response}}\ \otimes\ \underbrace{\Delta_q(D_1,H_1^{\sphericalangle})}_{\text{DiFF}}$};
\draw[line] (1.5,8.69)--(26.5,8.69);
\node[anchor=west,font=\fontsize{20}{28}\selectfont] at (2.0,7.57) {\textcolor{teal}{$\bullet$}\quad 识别 hard、collinear、soft 动量区域};
\node[anchor=west,font=\fontsize{20}{28}\selectfont] at (2.0,5.94) {\textcolor{teal}{$\bullet$}\quad PDF / DiFF 的算符定义与 Wilson lines};
\node[anchor=west,font=\fontsize{20}{28}\selectfont] at (2.0,4.31) {\textcolor{teal}{$\bullet$}\quad 自旋投影与 hard coefficient 的匹配};
\node[anchor=west,font=\fontsize{20}{28}\selectfont] at (2.0,2.68) {\textcolor{teal}{$\bullet$}\quad collinear singularities 的吸收与尺度演化};
\finishslide

\slide{7}{从这个课题学习 helicity amplitudes}{固定 $\phi_e=0$；[5] Sec. 8.1、[6] Sec. 2}
\node[anchor=west,font=\fontsize{20}{28}\selectfont] at (1.2,11.67) {$\rho^e=\frac12\begin{pmatrix}1&P_{eT}\\P_{eT}&1\end{pmatrix}$};
\node[anchor=west,font=\fontsize{21}{29}\selectfont] at (9.6,11.67) {$W^q_{\lambda\lambda'}=\frac12\sum_{a,a',b,c}\rho^e_{aa'}\,\mathcal M_{c\lambda;ab}\,\mathcal M^*_{c\lambda';a'b}$};
\node[anchor=west,font=\smallfont,text=gray] at (1.2,9.85) {$a,a':e_{\rm in}\qquad b:q_{\rm in}\qquad c:e_{\rm out}\qquad \lambda,\lambda':q_{\rm out}$};
\draw[line] (1.2,8.81)--(26.8,8.81);
\node[anchor=west,font=\fontsize{20}{28}\selectfont] at (2,7.59) {\textcolor{teal}{$\bullet$}\quad 构造 spinors、helicity projectors 与相位 convention};
\node[anchor=west,font=\fontsize{20}{28}\selectfont] at (2,5.96) {\textcolor{teal}{$\bullet$}\quad 计算 SM--scalar/tensor 的线性干涉};
\node[anchor=west,font=\fontsize{20}{28}\selectfont] at (2,4.33) {\textcolor{teal}{$\bullet$}\quad 从振幅构造末态 quark 自旋密度矩阵 $W$};
\node[anchor=west,font=\fontsize{20}{28}\selectfont] at (2,2.70) {\textcolor{teal}{$\bullet$}\quad 提取 $W_{+-}$，与 DiFF 收缩得到角分布};
\finishslide

\slide{8}{可观测量：双强子方位角 $\phi_R$}{固定 $\phi_e=0$：电子横自旋相对散射面固定；图为角分布示意}
\node[anchor=west,font=\fontsize{22}{30}\selectfont] at (1.2,11.80) {$\left.\frac{\dd\sigma^\pm}{\dd\phi_R}\right|_{\phi_e=0}\propto1\pm P_{eT}\bigl[a\sin\phi_R+b\cos\phi_R\bigr]$};
\node[anchor=west,font=\fontsize{20}{28}\selectfont] at (1.2,8.82) {$A^w=\dfrac{\int_0^{2\pi}w(\phi_R)(\dd\sigma^+-\dd\sigma^-)}{P_{eT}\int_0^{2\pi}(\dd\sigma^++\dd\sigma^-)}$};
\node[anchor=west,font=\fontsize{22}{30}\selectfont,text=blue] at (1.2,5.72) {$A^{\sin\phi_R}=\frac a2,\qquad A^{\cos\phi_R}=\frac b2$};
\node[anchor=west,font=\bodyfont,text=gray] at (1.2,3.71) {$+/-$：翻转电子横自旋};
\begin{scope}[shift={(16.5,6.1)},x=1.35cm,y=2cm]
\draw[->,gray] (0,-.95)--(0,1.13) node[above,font=\smallfont] {归一化角分布};
\draw[->,gray] (-.1,-.85)--(6.6,-.85) node[right,font=\smallfont] {$\phi_R$};
\draw[dashed,gray] (0,0)--(6.28,0);
\draw[blue,line width=1.7pt,domain=0:6.283,samples=80] plot(\x,{.6*sin(\x r)});
\draw[orange,line width=1.7pt,domain=0:6.283,samples=80] plot(\x,{-.6*sin(\x r)});
\node[anchor=west,font=\smallfont,text=blue] at (2.8,1.00) {$+S_{eT}$};
\node[anchor=west,font=\smallfont,text=orange] at (4.8,1.00) {$-S_{eT}$};
\node[anchor=east,font=\smallfont] at (-.13,0) {$1$};
\node[below,font=\smallfont] at (0,-.91) {$0$};
\node[below,font=\smallfont] at (3.14,-.91) {$\pi$};
\node[below,font=\smallfont] at (6.28,-.91) {$2\pi$};
\end{scope}
\node[font=\smallfont,text=gray] at (21,3.28) {示意图取 $b=0$};
\finishslide


\slide{9}{形成论文所需的结果}{研究计划；reach / uncertainty 曲线待计算}
\txt{1}{12.64}{25.6}{预期的科学结果是：建立四费米子干涉的末态 DiFF observable，并量化它与原 $h_1$ 方案的互补性。是否给出更强约束，要在相同实验假设下比较。}
\draw[line,line width=2pt] (2.0,9.72)--(25.8,9.72);
\foreach \xx/\nn in {2/1,10.1/2,18.6/3}{\fill[teal] (\xx,9.72) circle (.19);\node[font=\smallfont,text=teal,above] at (\xx,10.13) {\nn};}
\txt{1}{9.02}{7.6}{\raggedright\textbf{解析结果}\par\vspace{7pt}复现 [1]；推导新 hard kernels、联合角权重及各算符的 \mbox{flavor 结构。}\par}
\txt{9.9}{9.02}{7.6}{\textbf{数值 forecast}\par\vspace{7pt}用 $\pi^+\pi^-$ 输入，给出含误差的角矩与耦合 contours；比较单参数和联合拟合。}
\txt{18.9}{9.02}{8.0}{\textbf{同条件比较}\par\vspace{7pt}统一总 luminosity 与 DIS cuts，纳入 pair efficiency，比较 $h_1$ 与 DiFF 两条路线。}
\draw[line] (1,5.28)--(27,5.28);
\sm{1}{4.84}{25.6}{\textbf{目标期刊：}\href{https://journals.aps.org/prd/}{Physical Review D（PRD）。}}
\finishslide

\slide{10}{知识准备与最短阅读路径}{参考文献链接可点击}
\txt{1}{12.68}{25.6}{阅读目标：能解释因子化中每个函数的物理含义；能独立计算并核查保留 helicity coherence 的 hard matrix。}
\node[anchor=west,font=\bodyfont\bfseries,text=teal] at (1,10.85) {物理基线与 DiFF 输入};
\sm{1}{10.15}{12.2}{\textbf{[1]} Wang, Wen, Xing, Yan (2024)\par
\src{2401.08419}{Four-fermion DSA at EIC / 2401.08419}\par
Sec. III：复现 scalar/tensor 干涉与原角结构。}
\sm{1}{8.0}{12.2}{\textbf{[2]} Wen, Yan, Yu, Yuan (2026)\par
\src{2408.07255}{DiFF and light-quark dipoles / 2408.07255}\par
Eqs. (2)--(4)：掌握 $W\Delta$、角分析与归一化。}
\sm{1}{5.85}{12.2}{\textbf{[3]} Cocuzza et al. (JAM), \src{2308.14857}{2308.14857}\par
$D_1,\HT$ 联合提取。正确调用参数化，以同一 replica 传播分子分母的相关误差。}
\sm{1}{3.70}{12.2}{\textbf{[4]} Pitonyak et al., \src{2305.11995}{2305.11995}\par
extended DiFF 定义与演化。对齐变量、measure、相对动量权重及 $dM_h/dM_h^2$ convention。}
\draw[line] (13.7,10.9)--(13.7,2.2);
\node[anchor=west,font=\bodyfont\bfseries,text=teal] at (14.7,10.85) {本地笔记中的理论准备};
\sm{14.7}{10.15}{12.2}{\textbf{[5]} Zhite Yu, PhD thesis, \src{2308.13080}{2308.13080}\par
Sec. 8.1：spin-density matrix。\par 写出相位一致的 $\rho^e$ 和 outgoing-quark $W_{+-}$。\par\vspace{7pt}
Secs. 2.4, 2.6.3, 2.7.1：region approximation、DIS factorization、LO PDF spin projection。}
\sm{14.7}{6.3}{12.2}{\textbf{[6]} Barone--Drago--Ratcliffe\par
\src{hep-ph/0104283}{Transverse polarisation / hep-ph/0104283}\par
Sec. 2；Secs. 6.7.1、6.8：spin 基础与双强子 SIDIS。复现标准 $h_1\HT$ 角分布。}
\sm{14.7}{3.65}{12.2}{\textbf{辅助：}Peskin \& Schroeder Sec. 5.2、Appendix A.2。以 explicit spinors 和 Dirac trace 两条路径核查同一振幅。}
\sm{1}{1.52}{25.8}{本地依据：\texttt{proposal.tex}、\texttt{transverse\_spin.md} 与 WhiteBoard；已发表结果、选择规则分析与待完成的 forecast 分别标明。}
\finishslide
\end{document}
